\documentclass[11pt]{article}
\usepackage{amsmath,amssymb,amsthm,hyperref}
\newtheorem{theorem}{Theorem}
\newtheorem{corollary}{Corollary}
\newcommand{\E}{\mathbb E}
\title{A precision-dependent lower bound for approximately fair dice}
\author{Research note; independently reviewed}
\date{30 September 2026}
\begin{document}\maketitle
Throughout, logarithms are natural, and $\log(1/0)=+\infty$.

\begin{theorem}\label{main}
There is an absolute $c>0$ with the following property. Let $m\ge2$,
$0\le\varepsilon\le1/2$, and let $p_1,\ldots,p_m:[K]\to[0,1]$ be
nondecreasing on rows with positive probability weights $w_q$.
Suppose every failure pattern $S\subseteq[m]$ satisfies
\[
 \frac{1-\varepsilon}{(m+1)\binom m{|S|}}
 \le\E_w\prod_{j\in S}(1-p_j)\prod_{i\notin S}p_i
 \le\frac{1+\varepsilon}{(m+1)\binom m{|S|}}.
\]
Then both endpoint row weights are at most
\[
 \frac1{c m\min\{\log m,\log(1/\varepsilon)\}}.
\]
In particular, for equal weights,
\[
 K\ge c m\min\{\log m,\log(1/\varepsilon)\}.
\]
\end{theorem}

\begin{corollary}
Suppose every strict order of $n\ge3$ independent equiprobable finite
dice with no cross-die ties has probability in $[(1-\varepsilon)/n!,(1+\varepsilon)/n!]$,
where $0\le\varepsilon\le1/2$. Every die then has at least
\[
 c n\min\{\log n,\log(1/\varepsilon)\}
\]
faces, after changing the absolute constant.
\end{corollary}
\begin{proof}
For any distinguished die and any subset of the other dice, the event
that precisely that subset lies above it is a disjoint union of
$|S|!(n-1-|S|)!$ strict orders. Summing their relative bounds gives
the hypothesis of Theorem~\ref{main} with $m=n-1$.
\end{proof}

\begin{proof}[Proof of Theorem~\ref{main}]
Put $A=\prod_jp_j$ and $\eta(q)=(m+1)w_qA(q)$.
This is a positive measure whose mass lies in $[1-\varepsilon,1+\varepsilon]$;
it is not normalized to a probability measure. On $A>0$ put
$r_j=(1-p_j)/p_j$ and $\lambda=\sum_jr_j$, assigning infinite odds to
zero entries when defining bad rows.
For $B_q(z)=\prod_j(p_j+z(1-p_j))$, the coefficient of $z^k$ in
$(m+1)\E_wB_q(z)$ lies in $[1-\varepsilon,1+\varepsilon]$.
For a fixed index $j$, the coefficient of $z^k$ in
\[
 (m+1)\E_w\left[(1-p_j)\prod_{i\ne j}(p_i+z(1-p_i))\right]
\]
lies between $(1-\varepsilon)(k+1)/m$ and
$(1+\varepsilon)(k+1)/m$.

Take $\tau=m^{-1/2}$, let $D=\{\max_jr_j\ge\tau\}$, and put $G=D^c$.
The bad set is an initial segment. If it is nonempty, choose one
column witnessing the inequality at its last row. On every bad row,
its success probability is at most $\tau^{-1}$ times its failure
probability. Coefficientwise comparison of the undivided polynomials,
which remains valid at zero entries, gives
\[
 1-\varepsilon-(1+\varepsilon)
             \frac{(k+1)/\tau+k}{m}
 \le a_k:=\sum_{q\in G}\eta(q)e_k(r(q))\le1+\varepsilon
 \qquad(0\le k<m).
\]
Writing
$M_k=\sum_{q\in G}\eta(q)\lambda(q)^k/k!$, the collision expansion
of ordered index tuples gives
\[
 0\le M_k-a_k\le\frac{(k-1)\tau}{2}M_{k-1}\qquad(k\ge2),
 \quad M_1=a_1.
\]
The same induction as in the exact endpoint lemma yields
$M_k\le(1+\varepsilon)(1+k\tau)$ for $k\tau\le1$.
Combining both estimates,
\begin{equation}\label{momenterror}
 |M_k-1|\le\varepsilon+\frac{3(k+1)}{\sqrt m}
 \qquad(0\le k\le\sqrt m,\ k<m).
\end{equation}
No probability renormalization was used in this argument.

We use the positive-measure Laguerre lemma proved in
\path{../../size_bounds/logarithmic_per_die_lower_bound.tex}:
if normalized moments through degree $2d-1$ have error at most $\delta$
and $\delta d9^d<1/6$, then some positive mass lies below $2/d$, and
any atom at $x\le2/d$ has mass at most
$32/d+1024\delta9^d/d^2$. The lemma explicitly permits a measure
whose total mass is not one.

Set
\[
 H=\min\{\log m,\log(1/\varepsilon)\},\qquad
 d=\left\lfloor\frac{H}{16\log3}\right\rfloor.
\]
For sufficiently large $H$, $d\ge2$, $2d-1\le\sqrt m$, and
\eqref{momenterror} gives $\delta\le\varepsilon+6d/\sqrt m$.
Since $m\ge e^H$, $\varepsilon\le e^{-H}$, and $9^d\le e^{H/8}$,
\[
 \delta d9^d\le d e^{-7H/8}+6d^2e^{-3H/8}\longrightarrow0.
\]
The corresponding atom error is $o(1/d)$ as well. Hence the last
row is good, has $\lambda(K)<2/d$, and its $\eta$-mass is at most
$33/d$, once $H$ exceeds an absolute constant. As
$A(K)\ge e^{-\lambda(K)}\ge e^{-2/d}$, this gives
$w_K\le C/(mH)$.

For completeness, the remaining bounded range of $H$ follows from a
uniform linear bound. If $m\ge36$, then
\[
 \eta(G)\ge1-\varepsilon-\frac{1+\varepsilon}{\sqrt m}\ge\frac14.
\]
At $z=1/2$, the coefficient bounds give
$(m+1)\E_w B_q(1/2)\le3$.
On $G$, $\log(1+r_j/2)\ge r_j/4$, so
\[
 \int_G e^{\lambda/4}\,d\eta\le3.
\]
Since $e^4/4>3$, some good row has $\lambda\le16$.
The last row therefore has $A(K)\ge e^{-16}$ and
\[
 (m+1)w_K e^{-16}\le\eta(K)\le1+\varepsilon\le3/2.
\]
This proves $w_K\le C/m$ with an absolute constant. When $m<36$,
the same estimate follows by increasing that constant, because
$w_K\le1$. This handles bounded $H$ and completes the last-row bound.
Finally reverse the row order and replace every $p_j$ by $1-p_j$.
The failure-pattern hypothesis is unchanged, proving the first-row bound.
\end{proof}

The theorem is a general lower bound for pointwise relative approximation.
For fixed nonzero error it is linear in $n$; it does not match the known
random-palindrome upper bound or establish a superlinear lower bound
at fixed accuracy. Its additional factor grows as the requested error
shrinks, up to the exact-fairness factor $\log n$.
\end{document}
